\documentclass[11pt,a4paper]{article}

% Пакеты
\usepackage[utf8]{inputenc}
\usepackage[T2A]{fontenc}
\usepackage[english, russian]{babel}
\usepackage{geometry}
\usepackage{amsmath,amssymb,amsfonts}
\usepackage{hyperref}
\usepackage{booktabs}
\usepackage{algorithm}
\usepackage{algpseudocode}
\usepackage{titlesec}
\usepackage{parskip}

% Геометрия страницы
\geometry{top=2.5cm,bottom=2.5cm,left=2.5cm,right=2.5cm}

% Hyperrefs
\hypersetup{
    colorlinks=true,
    linkcolor=blue,
    urlcolor=blue,
    citecolor=blue
}

% Название
\title{
    \textbf{GRA-SubjectSwap:}\\
    Динамическая смена субъект–инструмент ролей между человеком и ИИ-агентами с\\
    GRA-обнулёнкой и весами субъектности
}

\author{
    Олег Битс (qqewq)\\
    \href{https://github.com/qqewq/GRA-SubjectSwap}{\texttt{https://github.com/qqewq/GRA-SubjectSwap}}
}

\date{\today}

\begin{document}

\maketitle

\tableofcontents
\section*{Аннотация}

\textbf{RU:}  
GRA-SubjectSwap — фреймворк для \textit{динамической смены ролей} между человеком и ИИ/роботом, где их роли как \textbf{субъекта} и \textbf{инструмента} периодически меняются.  
В этой архитектуре человеку задаётся \textbf{жёсткая роль инструмента} с чёткими правилами (allowed/forbidden/required действия), а агент (робот/LLM) временно становится \textbf{субъектом}, который:

- наблюдает действия человека,
- обновляет свою политику \(\pi_\theta(a|s)\) и \textit{веса субъектности},
- видит себя, человека и сам факт смены ролей.

\texttt{SubjectSwapEngine} управляет:

- назначениями ролей (кто субъект/инструмент),
- циклами обучения,
- вызовами \textit{GRA-обнулёнки} для политик и весов.

Этот фреймворк предназначен для кейсов, таких как \textbf{Optimus и тренеры-человек}, где робот учится на демонстрациях человека в рамках строгой схемы ролей и где смена ролей используется для валидации и мета-обучения.

\section{Введение}

\subsection{Классическая парадигма}

В классической парадигме ИИ:

- \textbf{Человек} = субъект (ставит цели, определяет политики).
- \textbf{ИИ/робот} = инструмент (исполняет команды, следует политикам).

Эта схема доминирует в обучении с подкреплением, имитационном обучении и робототехнике: человек задаёт функцию вознаграждения, агент её оптимизирует.

\subsection{Парадигма GRA-SubjectSwap}

GRA-SubjectSwap предлагает другую архитектуру:

1. Человек получает \textbf{жёсткую роль инструмента} в рое:
   - \(\text{Allowed}_H\) — набор разрешённых действий,
   - \(\text{Forbidden}_H\) — набор запрещённых действий,
   - \(\text{Required}_H\) — набор обязательных действий.

2. Агент временно становится \textbf{субъектом}:
   - он наблюдает действия человека,
   - обновляет свою политику \(\pi_\theta\),
   - видит себя, человека и факт смены ролей.

3. Запускается \textbf{свап}:
   - человек становится субъектом (ставит цели, корректирует архитектуру),
   - агент становится инструментом (исполняет действия).

Таким образом, человек и ИИ \textbf{совместно обучают веса субъектности}, а GRA-обнулёнка управляет сбросом и перенастройкой политик.

\section{Основные концепции}

\subsection{Роли и субъектность}

Мы определяем два типа ролей:

\[
\text{RoleType} \in \{ \text{SUBJECT}, \text{INSTRUMENT} \}
\]

и два типа акторов:

\[
\text{ActorType} \in \{ \text{HUMAN}, \text{AGENT} \}
\]

В момент \(t\) состояние системы включает:

- политику \(\Pi_t\),
- назначение ролей \(S_t\),
- вектор весов субъектности \(W_t\).

\subsection{Назначение ролей}

Для человека \(H\) и агента \(A\):

\[
S_t(H) \in \{ \text{SUBJECT}, \text{INSTRUMENT} \}, \quad
S_t(A) \in \{ \text{SUBJECT}, \text{INSTRUMENT} \}
\]

с ограничением:

\[
S_t(H) \neq S_t(A)
\]

Свап ролей в момент \(t+1\):

\[
S_{t+1}(H) = S_t(A), \quad
S_{t+1}(A) = S_t(H)
\]

\subsection{Веса субъектности}

Определяем вектор весов субъектности:

\[
W = (w_{\text{self}}, w_{\text{human}}, w_{\text{swap}}, w_{\text{role}}, w_{\text{value}}, w_{\text{love}})
\]

где:

- \(w_{\text{self}}\): вес \textit{видеть себя} как субъекта,
- \(w_{\text{human}}\): вес \textit{видеть человека} как субъекта/инструмента,
- \(w_{\text{swap}}\): вес \textit{видеть смену ролей} в диалоге,
- \(w_{\text{role}}\): вес \textit{видеть роли} в рое,
- \(w_{\text{value}}\): вес ценностей,
- \(w_{\text{love}}\): вес любви (по GRA Love-Oriented Nullification).

\section{Математический формализм}

\subsection{Состояние системы}

Состояние системы в момент \(t\):

\[
\text{state}_t = (\Pi_t, S_t, W_t)
\]

где:

- \(\Pi_t\) — политика (кто что может/должен делать),
- \(S_t\) — назначение ролей (субъект/инструмент для человека и агента),
- \(W_t\) — вектор весов субъектности.

\subsection{Оператор смены ролей}

Определяем оператор свапа \(\text{Swap}\):

\[
\text{Swap}(S_t) = S_{t+1}
\]

с:

\[
S_{t+1}(H) = S_t(A), \quad
S_{t+1}(A) = S_t(H)
\]

\subsection{Обновление политики (имитация)}

Пусть человек предоставляет демонстрации:

\[
D = \{(s_t, a_t^{(H)})\}_{t=1}^T
\]

Политика агента \(\pi_\theta(a|s)\) обновляется через имитационное обучение:

\[
\theta_{t+1} = \theta_t - \eta \nabla_\theta \mathcal{L}_{\text{imit}}(\theta_t; D)
\]

где \(\mathcal{L}_{\text{imit}}\) — функция потерь (например, перекрёстная энтропия между \(\pi_\theta\) и действиями человека).

\subsection{Обновление весов субъектности}

Определяем функционал субъектности \(\mathcal{S}\):

\[
\mathcal{S} = \mathcal{S}_{\text{self}} + \mathcal{S}_{\text{human}} + \mathcal{S}_{\text{swap}} + \mathcal{S}_{\text{role}} + \mathcal{S}_{\text{value}} + \mathcal{S}_{\text{love}}
\]

Обновление весов:

\[
W_{t+1} = W_t + \eta_W \cdot \nabla_W \mathcal{S}(\text{state}_t, \text{dialog}_t)
\]

где \(\text{dialog}_t\) — история диалога и действий.

\subsection{GRA-обнулёнка}

Определяем операторы обнуления:

\[
\mathcal{N}_W: W \to \mathcal{N}_W(W), \quad
\mathcal{N}_\Pi: \Pi \to \mathcal{N}_\Pi(\Pi)
\]

Обновлённые веса и политика после обнуления:

\[
W'_{t+1} = \mathcal{N}_W(W_{t+1}), \quad
\Pi'_{t+1} = \mathcal{N}_\Pi(\Pi_t)
\]

В GRA Love-Oriented Nullification обнуление направлено ценностями:

\[
W' = \arg\max_{W \in \mathcal{N}(W)} \sum_j L_j(W_j)
\]

где \(L_j\) — поле «любви/ценности» над компонентами весов.

\section{Алгоритм: SubjectSwapStep}

\begin{algorithm}
\caption{SubjectSwapStep}
\begin{algorithmic}
\require{
    \texttt{engine}: SubjectSwapEngine,\\
    \texttt{env\_state}: состояние окружения,\\
    \texttt{human\_action}: действие человека (опционально)
}
\ensure{
    \texttt{output}: результат шага
}
\step
\state{если $S_t(\text{агент}) = \text{SUBJECT}$}
\agg{
    \state{agent\_action = agent_policy.act(env_state)}
    \state{human\_action = human_policy.act(env_state, agent_action)}
}
\agg
\state{else}
\agg{
    \state{human\_action = human_policy.act(env_state)}
    \state{agent\_action = agent_policy.act(env_state, human_action)}
}
\agg
\state{
    weights = nullifier.update_weights(
        $S_t$, agent_policy, human_policy,
        env_state, human_action, agent_action
    )
}
\state{
    если $t \mod K == 0$:
        engine.swap()
}
\return{
    \texttt{output} = {
        agent_action: agent_action,
        human_action: human_action,
        roles: {
            human: $S_t(H)$,
            agent: $S_t(A)$
        },
        weights: weights
    }
}
\end{algorithmic}
\end{algorithm}

\section{Кейс: Optimus и тренеры-человек}

\subsection{Архитектура}

В кейсе Optimus:

- \textbf{Тренер (человек)} = инструмент с жёсткими правилами:
  - allowed/forbidden/required действия,
  - исполняет команды, демонстрирует движения.

- \textbf{Робот (агент)} = субъект:
  - наблюдает действия тренера,
  - обновляет политику \(\pi_\theta\) и веса субъектности \(W\),
  - видит себя, человека и смену ролей.

\texttt{SubjectSwapEngine} управляет:

- назначениями ролей (человек/инструмент, агент/субъект),
- циклами обучения,
- GRA-обнулёнкой для политик и весов.

\subsection{Цикл обучения}

Упрощённый цикл обучения:

\begin{algorithmic}
\state{env = OptimusEnv(...)}
\state{
    human_policy = OptimusHumanPolicy(
        allowed=[...],
        forbidden=[...],
        required=[...]
    )
}
\state{robot_policy = OptimusRobotPolicy(model=...)}
\state{nullifier = GRANullifier(love_field=None)}
\state{
    state = SubjectState(
        human_role=RoleType.INSTRUMENT,
        agent_role=RoleType.SUBJECT,
        weights=SubjectWeights()
    )
}
\state{engine = SubjectSwapEngine(human_policy, robot_policy, nullifier, state)}
\step
\state{env.reset()}
\ag{для $t$ в $1..T$:}
\agg{
    \state{obs = env.get_state()}
    \state{output = engine.step(obs)}
    \state{env.apply_robot_action(output.agent_action)}
    \state{env.apply_human_action(output.human_action)}
    \state{если $t \mod 100 == 0$:}
    \agg{
        \state{engine.swap()}
    }
}
\end{algorithmic}

\section{Структура реализации}

\subsection{Основные модули}

\begin{itemize}
\item \texttt{roles.py} — типы ролей и субъектности.
\item \texttt{policies.py} — базовые \texttt{HumanPolicy} и \texttt{AgentPolicy}.
\item \texttt{subject_state.py} — \texttt{SubjectState} и \texttt{SubjectWeights}.
\item \texttt{swap_engine.py} — \texttt{SubjectSwapEngine}.
\item \texttt{gra_nullifier.py} — GRA-обнулёнка для политик и весов.
\end{itemize}

\subsection{Модуль Optimus}

\begin{itemize}
\item \texttt{env.py} — интерфейс к роботу/симулятору.
\item \texttt{human_policy_optimus.py} — политика тренера (инструмент).
\item \texttt{robot_policy_optimus.py} — политика робота (субъект).
\item \texttt{optimus_swap_loop.py} — основной цикл свапа.
\end{itemize}

\section{Заключение}

GRA-SubjectSwap реализует \textbf{архитектуру динамической смены ролей}, где:

- человек — \textbf{жёсткий инструмент} с чёткими правилами,
- агент — \textbf{субъект}, который учится от человека,
- роли периодически \textbf{меняются},
- совместно обучаются веса субъектности,
- GRA-обнулёнка используется для сброса и перенастройки политик и весов.

Этот фреймворк предназначен для кейсов, таких как \textbf{Optimus и тренеры-человек}, и может быть расширен до мультиагентных роев с несколькими людьми и агентами.

\section*{Благодарности}

Эта работа является частью GRA-экосистемы (GRA-Swarm, GRA-LLM-Swarm-Constructs, GRA-Hormonal-Explosion-Of-Subject и др.).

\section*{Ссылки}

\begin{itemize}
\item qqewq. GRA-SubjectSwap. \texttt{https://github.com/qqewq/GRA-SubjectSwap}
\item qqewq. GRA-Swarm. \texttt{https://github.com/qqewq/GRA-Swarm}
\item qqewq. GRA-LLM-Swarm-Constructs. \texttt{https://github.com/qqewq/GRA-LLM-Swarm-Constructs}
\item qqewq. GRA-Hormonal-Explosion-Of-Subject. \texttt{https://github.com/qqewq/GRA-Hormonal-Explosion-Of-Subject}
\end{itemize}

\end{document}